\documentclass[10pt,a4paper]{amsart}
\usepackage{amsmath,amssymb,amsthm,graphicx}
% Retained minimal preamble for the Stacks Project chapter more-algebra.tex (More on Algebra).
% The chapter file pinned-original-source.tex carries only \input{preamble} on its first line, so no
% preamble text lives inside the chapter itself.  The Stacks Project uses one single shared preamble for
% every chapter; the snapshot stages it identically as sources/stacks-chow/preamble.tex and
% sources/stacks-etale-cohomology/preamble.tex.  The block below is copied mechanically, line for line,
% from sources/stacks-chow/preamble.tex lines 14-18 (the Xy-pic setup) and 154-217 (the native theorem
% declarations together with the mathematical macros).  Nothing was added, removed or redefined; the
% wrapper class and the amsmath/amssymb/amsthm/graphicx base packages are supplied by the assembler.
% Chapter: more-algebra.tex -- More on Algebra.
% For commutative diagrams we use Xy-pic
\usepackage[all]{xy}

% We use 2cell for 2-commutative diagrams.
\xyoption{2cell}

% Theorem environments.
%
\theoremstyle{plain}
\newtheorem{theorem}[subsection]{Theorem}
\newtheorem{proposition}[subsection]{Proposition}
\newtheorem{lemma}[subsection]{Lemma}

\theoremstyle{definition}
\newtheorem{definition}[subsection]{Definition}
\newtheorem{example}[subsection]{Example}
\newtheorem{exercise}[subsection]{Exercise}
\newtheorem{situation}[subsection]{Situation}

\theoremstyle{remark}
\newtheorem{remark}[subsection]{Remark}
\newtheorem{remarks}[subsection]{Remarks}

\numberwithin{equation}{subsection}

% Macros
%
\def\lim{\mathop{\mathrm{lim}}\nolimits}
\def\colim{\mathop{\mathrm{colim}}\nolimits}
\def\Spec{\mathop{\mathrm{Spec}}}
\def\Hom{\mathop{\mathrm{Hom}}\nolimits}
\def\Ext{\mathop{\mathrm{Ext}}\nolimits}
\def\SheafHom{\mathop{\mathcal{H}\!\mathit{om}}\nolimits}
\def\SheafExt{\mathop{\mathcal{E}\!\mathit{xt}}\nolimits}
\def\Sch{\mathit{Sch}}
\def\Mor{\mathop{\mathrm{Mor}}\nolimits}
\def\Ob{\mathop{\mathrm{Ob}}\nolimits}
\def\Sh{\mathop{\mathit{Sh}}\nolimits}
\def\NL{\mathop{N\!L}\nolimits}
\def\CH{\mathop{\mathrm{CH}}\nolimits}
\def\proetale{{pro\text{-}\acute{e}tale}}
\def\etale{{\acute{e}tale}}
\def\QCoh{\mathit{QCoh}}
\def\Ker{\mathop{\mathrm{Ker}}}
\def\Im{\mathop{\mathrm{Im}}}
\def\Coker{\mathop{\mathrm{Coker}}}
\def\Coim{\mathop{\mathrm{Coim}}}

% Boxtimes
%
\DeclareMathSymbol{\boxtimes}{\mathbin}{AMSa}{"02}

%
% Macros for moduli stacks/spaces
%
\def\QCohstack{\mathcal{QC}\!\mathit{oh}}
\def\Cohstack{\mathcal{C}\!\mathit{oh}}
\def\Spacesstack{\mathcal{S}\!\mathit{paces}}
\def\Quotfunctor{\mathrm{Quot}}
\def\Hilbfunctor{\mathrm{Hilb}}
\def\Curvesstack{\mathcal{C}\!\mathit{urves}}
\def\Polarizedstack{\mathcal{P}\!\mathit{olarized}}
\def\Complexesstack{\mathcal{C}\!\mathit{omplexes}}
% \Pic is the operator that assigns to X its picard group, usage \Pic(X)
% \Picardstack_{X/B} denotes the Picard stack of X over B
% \Picardfunctor_{X/B} denotes the Picard functor of X over B
\def\Pic{\mathop{\mathrm{Pic}}\nolimits}
\def\Picardstack{\mathcal{P}\!\mathit{ic}}
\def\Picardfunctor{\mathrm{Pic}}
\def\Deformationcategory{\mathcal{D}\!\mathit{ef}}
\usepackage{cleveref}
\crefname{originallemma}{original lemma}{original lemmas}
\Crefname{originallemma}{Original lemma}{Original lemmas}
\makeatletter
\newlabel{more-algebra-lemma-double-base-change}{{08YU}{source}{}{sourcepointer-1}{}}
\newlabel{more-algebra-lemma-double-base-change@cref}{{[originallemma][0][]08YU}{[1][0][]source}{}{}{}}
\makeatother
\crefname{originallemma}{original lemma}{original lemmas}
\Crefname{originallemma}{Original lemma}{Original lemmas}
\makeatletter
\newlabel{lemma-K-flat-resolution}{{06Y4}{source}{}{sourcepointer-2}{}}
\newlabel{lemma-K-flat-resolution@cref}{{[originallemma][0][]06Y4}{[1][0][]source}{}{}{}}
\makeatother
\crefname{originallemma}{original lemma}{original lemmas}
\Crefname{originallemma}{Original lemma}{Original lemmas}
\makeatletter
\newlabel{lemma-functoriality-derived-base-change}{{0BYK}{source}{}{sourcepointer-3}{}}
\newlabel{lemma-functoriality-derived-base-change@cref}{{[originallemma][0][]0BYK}{[1][0][]source}{}{}{}}
\makeatother
\crefname{originallemma}{original lemma}{original lemmas}
\Crefname{originallemma}{Original lemma}{Original lemmas}
\makeatletter
\newlabel{lemma-base-change-K-flat}{{06Y1}{source}{}{sourcepointer-4}{}}
\newlabel{lemma-base-change-K-flat@cref}{{[originallemma][0][]06Y1}{[1][0][]source}{}{}{}}
\makeatother
\crefname{originallemma}{original lemma}{original lemmas}
\Crefname{originallemma}{Original lemma}{Original lemmas}
\makeatletter
\newlabel{lemma-tensor-product-K-flat}{{0795}{source}{}{sourcepointer-5}{}}
\newlabel{lemma-tensor-product-K-flat@cref}{{[originallemma][0][]0795}{[1][0][]source}{}{}{}}
\makeatother
\crefname{originallemma}{original lemma}{original lemmas}
\Crefname{originallemma}{Original lemma}{Original lemmas}
\makeatletter
\newlabel{homology-remark-triple-complex}{{08BI}{source}{}{sourcepointer-6}{}}
\newlabel{homology-remark-triple-complex@cref}{{[originallemma][0][]08BI}{[1][0][]source}{}{}{}}
\makeatother
\crefname{originallemma}{original lemma}{original lemmas}
\Crefname{originallemma}{Original lemma}{Original lemmas}
\makeatletter
\newlabel{lemma-derived-base-change}{{06Y6}{source}{}{sourcepointer-7}{}}
\newlabel{lemma-derived-base-change@cref}{{[originallemma][0][]06Y6}{[1][0][]source}{}{}{}}
\makeatother
\crefname{originallemma}{original lemma}{original lemmas}
\Crefname{originallemma}{Original lemma}{Original lemmas}
\makeatletter
\newlabel{lemma-flip-douoble-tensor-product}{{0BYI}{source}{}{sourcepointer-8}{}}
\newlabel{lemma-flip-douoble-tensor-product@cref}{{[originallemma][0][]0BYI}{[1][0][]source}{}{}{}}
\makeatother
\crefname{originallemma}{original lemma}{original lemmas}
\Crefname{originallemma}{Original lemma}{Original lemmas}
\makeatletter
\newlabel{lemma-triple-tensor-product}{{0BYJ}{source}{}{sourcepointer-9}{}}
\newlabel{lemma-triple-tensor-product@cref}{{[originallemma][0][]0BYJ}{[1][0][]source}{}{}{}}
\makeatother
\title{The two-fold derived base change comparison for a complex}
\author{The Stacks Project Authors}
\date{stacks-project:more-algebra.tex:local-snapshot-20261012; GNU FDL 1.2 or later}
\begin{document}
\maketitle
\setcounter{subsection}{1}
\noindent\textit{Editorial note.} All delivered mathematical text is the author's own, copied from the pinned Stacks Project chapter file more-algebra.tex; no statement or proof step was rewritten, reordered or removed.\par\medskip
\noindent\textit{Editorial note.} The chapter file pinned-original-source.tex carries only \textbackslash{}input\{preamble\} on its first line, so no preamble text lives inside the chapter. The retained file sources/stacks-more-algebra/retained-preamble.tex is a mechanical, line-for-line copy of the project's single shared preamble as staged in the snapshot (Xy-pic setup plus theorem declarations and mathematical macros). Nothing was added, removed or redefined, and the excerpt serves only as preamble, never as a span coordinate.\par\medskip
\noindent\textit{Editorial note.} Both counted spans start at the \textbackslash{}begin line of their environment, so the statement and the author's proof are delivered whole; surrounding prose of the section is deliberately outside the counted spans and is carried as standard prerequisites instead.\par\medskip
\noindent\textit{Editorial note.} The citation fields (source\_id, source\_version, authors, source\_url, license.name, source\_title) are identical to those of every other card this lane has produced from this chapter, so the batch selector treats them as one source.\par\medskip
\noindent\textit{Allowed standard prerequisites (editorial).}\begin{itemize}
\item K-flat complex, defined by the property that tensoring with it preserves quasi-isomorphisms. definition at source lines 14185-14191 of the same chapter. Supplies the kind of resolution used in the proof.
\item Derived tensor product may be computed with a K-flat representative. used at proof lines 14821-14860. Justifies replacing the derived constructions by ordinary ones.
\item Associativity of the total tensor product of complexes. used at proof lines 14860-14891. Performs the final identification of the two sides.
\end{itemize}\medskip
\ifdefined\hypertarget\hypertarget{sourcepointer-1}{}\fi
\noindent\textit{Source reference.} Original lemma 08YU: focused-seeds/AGP/sources/stacks-snapshot-20261011/tags.txt, source lines 11352--11352. The lemma's own canonical tag, quoted so that readers can resolve the label inside the delivered text; it points at the counted statement itself and carries no proof content.\par\medskip
\ifdefined\hypertarget\hypertarget{sourcepointer-2}{}\fi
\noindent\textit{Source reference.} Original lemma 06Y4: focused-seeds/AGP/sources/stacks-snapshot-20261011/tags.txt, source lines 8873--8873. Navigation pointer to a statement the counted proof refers to but does not contain; nothing load-bearing is delegated to it.\par\medskip
\ifdefined\hypertarget\hypertarget{sourcepointer-3}{}\fi
\noindent\textit{Source reference.} Original lemma 0BYK: focused-seeds/AGP/sources/stacks-snapshot-20261011/tags.txt, source lines 15018--15018. Navigation pointer to a statement the counted proof refers to but does not contain; nothing load-bearing is delegated to it.\par\medskip
\ifdefined\hypertarget\hypertarget{sourcepointer-4}{}\fi
\noindent\textit{Source reference.} Original lemma 06Y1: focused-seeds/AGP/sources/stacks-snapshot-20261011/tags.txt, source lines 8870--8870. Navigation pointer to a statement the counted proof refers to but does not contain; nothing load-bearing is delegated to it.\par\medskip
\ifdefined\hypertarget\hypertarget{sourcepointer-5}{}\fi
\noindent\textit{Source reference.} Original lemma 0795: focused-seeds/AGP/sources/stacks-snapshot-20261011/tags.txt, source lines 9260--9260. Navigation pointer to a statement the counted proof refers to but does not contain; nothing load-bearing is delegated to it.\par\medskip
\ifdefined\hypertarget\hypertarget{sourcepointer-6}{}\fi
\noindent\textit{Source reference.} Original lemma 08BI: focused-seeds/AGP/sources/stacks-snapshot-20261011/tags.txt, source lines 10568--10568. Navigation pointer to a statement the counted proof refers to but does not contain; nothing load-bearing is delegated to it.\par\medskip
\ifdefined\hypertarget\hypertarget{sourcepointer-7}{}\fi
\noindent\textit{Source reference.} Original lemma 06Y6: focused-seeds/AGP/sources/stacks-snapshot-20261011/tags.txt, source lines 8875--8875. Navigation pointer to a statement the counted proof refers to but does not contain; nothing load-bearing is delegated to it.\par\medskip
\ifdefined\hypertarget\hypertarget{sourcepointer-8}{}\fi
\noindent\textit{Source reference.} Original lemma 0BYI: focused-seeds/AGP/sources/stacks-snapshot-20261011/tags.txt, source lines 15016--15016. Navigation pointer to a statement the counted proof refers to but does not contain; nothing load-bearing is delegated to it.\par\medskip
\ifdefined\hypertarget\hypertarget{sourcepointer-9}{}\fi
\noindent\textit{Source reference.} Original lemma 0BYJ: focused-seeds/AGP/sources/stacks-snapshot-20261011/tags.txt, source lines 15017--15017. Navigation pointer to a statement the counted proof refers to but does not contain; nothing load-bearing is delegated to it.\par\medskip
\begin{lemma}
\label{lemma-double-base-change}
Let $A \to B \to C$ be ring maps. Let $N^\bullet$ be a complex of
$B$-modules and $K^\bullet$ a complex of $C$-modules.
The compositions of the functors
$$
D(A) \xrightarrow{- \otimes_A^\mathbf{L} N^\bullet}
D(B) \xrightarrow{- \otimes_B^\mathbf{L} K^\bullet} D(C)
$$
is the functor
$- \otimes_A^\mathbf{L} (N^\bullet \otimes_B^\mathbf{L} K^\bullet) :
D(A) \to D(C)$. If $M$, $N$, $K$ are modules over $A$, $B$, $C$, then
we have
$$
(M \otimes_A^\mathbf{L} N) \otimes_B^\mathbf{L} K =
M \otimes_A^\mathbf{L} (N \otimes_B^\mathbf{L} K) =
(M \otimes_A^\mathbf{L} C) \otimes_C^\mathbf{L} (N \otimes_B^\mathbf{L} K)
$$
in $D(C)$. We also have a canonical isomorphism
$$
(M \otimes_A^\mathbf{L} N) \otimes_B^\mathbf{L} K \longrightarrow
(M \otimes_A^\mathbf{L} K) \otimes_C^\mathbf{L} (N \otimes_B^\mathbf{L} C)
$$
using signs. Similar results holds for complexes.
\end{lemma}

\begin{proof}
Choose a K-flat complex $P^\bullet$ of $B$-modules and a quasi-isomorphism
$P^\bullet \to N^\bullet$ (Lemma \ref{lemma-K-flat-resolution}).
Let $M^\bullet$ be a K-flat complex of $A$-modules representing an
arbitrary object of $D(A)$. Then we see that
$$
(M^\bullet \otimes_A^\mathbf{L} P^\bullet) \otimes_B^\mathbf{L} K^\bullet
\longrightarrow
(M^\bullet \otimes_A^\mathbf{L} N^\bullet) \otimes_B^\mathbf{L} K^\bullet
$$
is an isomorphism by Lemma \ref{lemma-functoriality-derived-base-change}
applied to the material inside the brackets. By
Lemmas \ref{lemma-base-change-K-flat} and \ref{lemma-tensor-product-K-flat}
the complex
$$
\text{Tot}(M^\bullet \otimes_A P^\bullet) =
\text{Tot}((M^\bullet \otimes_R A) \otimes_A P^\bullet
$$
is K-flat as a complex of $B$-modules and it represents
the derived tensor product in $D(B)$ by construction.
Hence we see that
$(M^\bullet \otimes_A^\mathbf{L} P^\bullet) \otimes_B^\mathbf{L} K^\bullet$
is represented by the complex
$$
\text{Tot}(\text{Tot}(M^\bullet \otimes_A P^\bullet)\otimes_B K^\bullet) =
\text{Tot}(M^\bullet \otimes_A \text{Tot}(P^\bullet \otimes_B K^\bullet))
$$
of $C$-modules. Equality by
Homology, Remark \ref{homology-remark-triple-complex}.
Going back the way we came we see that this is equal to
$$
M^\bullet \otimes_A^\mathbf{L} (P^\bullet \otimes_B^\mathbf{L} K^\bullet)
\longleftarrow
M^\bullet \otimes_A^\mathbf{L} (N^\bullet \otimes_B^\mathbf{L} K^\bullet)
$$
The arrow is an isomorphism by definition of the functor
$-\otimes_B^\mathbf{L} K^\bullet$. All of these constructions
are functorial in the complex $M^\bullet$ and hence we obtain
our isomorphism of functors.

\medskip\noindent
By the above we have the first equality in
$$
(M \otimes_A^\mathbf{L} N) \otimes_B^\mathbf{L} K =
M \otimes_A^\mathbf{L} (N \otimes_B^\mathbf{L} K) =
(M \otimes_A^\mathbf{L} C) \otimes_C^\mathbf{L} (N \otimes_B^\mathbf{L} K)
$$
The second equality follows from the final statement of
Lemma \ref{lemma-derived-base-change}.
The same thing allows us to write
$N \otimes_B^\mathbf{L} K = (N \otimes_B^\mathbf{L} C) \otimes_C^\mathbf{L} K$
and substituting we get
\begin{align*}
(M \otimes_A^\mathbf{L} N) \otimes_B^\mathbf{L} K
& =
(M \otimes_A^\mathbf{L} C) \otimes_C^\mathbf{L}
((N \otimes_B^\mathbf{L} C) \otimes_C^\mathbf{L} K) \\
& =
(M \otimes_A^\mathbf{L} C) \otimes_C^\mathbf{L}
(K \otimes_C^\mathbf{L} (N \otimes_B^\mathbf{L} C)) \\
& =
((M \otimes_A^\mathbf{L} C) \otimes_C^\mathbf{L} K)
\otimes_C^\mathbf{L} (N \otimes_B^\mathbf{L} C)) \\
& =
(M \otimes_C^\mathbf{L} K)
\otimes_C^\mathbf{L} (N \otimes_B^\mathbf{L} C)
\end{align*}
by Lemmas \ref{lemma-flip-douoble-tensor-product} and
\ref{lemma-triple-tensor-product}
as well as the previously mentioned lemma.
\end{proof}

\end{document}
